\section{Full coefficient space and open realization}
\label{sec:quadratic-transfer}
Identify the space $\mathcal Q_2$ of real planar quadratic vector
fields with $\R^{12}$ by their coefficients. The affine and positive-time
normalization of Dumortier and Rousseau transfers the whole-graphic
count to this full coefficient space.

\begin{corollary}\label{cor:quadratic-space}
Fix $b_0\in(0,2)$. There are an open coefficient neighborhood
$\mathcal W\subset\mathcal Q_2$ of $X_{b_0,0}$ and $\eta>0$ such
that every $Q\in\mathcal W$ satisfies
\[
 \sum_{\substack{L\ {\rm isolated\ periodic\ orbit\ of}\ Q\\
                    d_{\rm H}(L,\Gamma_{b_0})<\eta}}m_Q(L)\le2.
\]
There is a sequence $Q_n\to X_{b_0,0}$ in $\mathcal Q_2$ for which
each $Q_n$ has two distinct isolated periodic orbits converging
separately to $\Gamma_{b_0}$ in two-sided Hausdorff distance.
Both orbits are eventually hyperbolic. Thus the whole-graphic
cyclicity of $\Gamma_{b_0}$ in the full quadratic coefficient space
is exactly two.
\end{corollary}

\begin{proof}
We verify the local normalization of
\cite[Proposition~2.1]{DR2009} by computing its differential.
For $d\in\R^2$, $A\in GL(2,\R)$ and $c>0$, set
\[
 \Phi(d,A,c,b,\mu)(u)
   =cA X_{b,\mu}\bigl(A^{-1}(u-d)\bigr).
\]
Its domain and target both have dimension twelve. At
$(0,I,1,b_0,0)$, an affine variation
$V=(p+az+\ell y,q+ez+fy)$ contributes $DVX-DXV$, where
$X=X_{b_0,0}$; write $\kappa,\beta,\nu$ for the variations of
$c,b,\mu$. If the total variation vanishes, its second component is
\[
 \nu_4-qz+(e-p)y-ez^2+(b_0e-a+\kappa)zy-(e+\ell)y^2=0.
\]
Hence $q=e=p=\ell=\nu_4=0$ and $\kappa=a$.
The $y$ and $y^2$ coefficients of the first component now give
$f=2a$ and $2a=0$. Its remaining coefficients give
$\beta=\nu_1=\nu_2=\nu_3=0$. The differential is therefore
invertible. The analytic inverse-function theorem supplies, for all
$Q$ in one coefficient neighborhood, analytic choices
$H_Q(u)=A_Qu+d_Q$, $c_Q>0$, $b_Q$ and $\mu_Q$ such that
\[
 Q(H_Q(u))=c_QA_QX_{b_Q,\mu_Q}(u),\qquad
 (H_Q,c_Q,b_Q,\mu_Q)\longrightarrow(\mathrm{id},1,b_0,0).
\]

Choose a compact interval $J\Subset(0,2)$ with $b_0$ in its interior,
and let $K,V,\rho_{\rm H}$ be given by
Theorem~\ref{thm:main}(i). The fast terminal intersection satisfies
$y'=f_b<0$, so continuous dependence and transversality give
Hausdorff continuity of $b\mapsto\Gamma_b$
(Appendix~\ref{sec:geometry}). Consequently
\[
 E_Q:=d_{\rm H}(H_Q\Gamma_{b_Q},\Gamma_{b_0})\longrightarrow0.
\]
After shrinking the coefficient neighborhood, let $M\ge1$ bound
both $\|A_Q\|$ and $\|A_Q^{-1}\|$. For $\widetilde L=H_Q^{-1}L$,
the two triangle inequalities give
\begin{align*}
 d_{\rm H}(\widetilde L,\Gamma_{b_Q})
 &\le M\bigl(d_{\rm H}(L,\Gamma_{b_0})+E_Q\bigr),\\
 d_{\rm H}(L,\Gamma_{b_0})
 &\le M d_{\rm H}(\widetilde L,\Gamma_{b_Q})+E_Q.
\end{align*}
Choose $\kappa_0>0$ so that the $\kappa_0$-neighborhood of
$\bigcup_{b\in J}\Gamma_b$ is contained in $K$, and set
$\theta=\min(\kappa_0,\rho_{\rm H})$.
Shrink $\mathcal W$ so that $b_Q\in J$, $\mu_Q\in V$ and
$E_Q<\theta/(4M)$, and take $\eta=\theta/(4M)$.
Every orbit counted in the corollary then gives
$\widetilde L\subset K$ and
$d_{\rm H}(\widetilde L,\Gamma_{b_Q})<\theta/2$.
These inequalities also show that fixed-base and moving-graphic
Hausdorff convergence are equivalent as $Q\to X_{b_0,0}$.

The positive constant time factor and affine map preserve isolated
periodic orbits, their minimal returns and their multiplicities, by
Section~\ref{sec:return-framework}. Theorem~\ref{thm:main}(i) gives
the upper bound.

Finally, the full coefficient space contains
$Q_n=X_{b_0,\mu_n}$ from Theorem~\ref{thm:main}(ii). These fields
eventually belong to $\mathcal W$ and realize two separate
whole-graphic limits at the fixed modulus $b_0$.
This proves sharpness in $\mathcal Q_2$.
\end{proof}

\begin{corollary}[Open realization]\label{cor:open-quadratic-realization}
Fix $b_0\in(0,2)$, and let $\mathcal W$ and $\eta$ be as in
Corollary~\ref{cor:quadratic-space}. There are fields
$Q_n\to X_{b_0,0}$ and nonempty open coefficient balls $U_n$
centered at $Q_n$, with
\[
 Q_n\in U_n\subset\mathcal W\cap B(Q_n,1/n)\subset\mathcal Q_2
\]
such that, for every $Q\in U_n$, the isolated periodic orbits satisfying
$d_{\rm H}(L,\Gamma_{b_0})<\eta$ consist exactly of two hyperbolic
cycles $L_n^-(Q)$ and $L_n^+(Q)$. Their characteristic multipliers
satisfy
\[
 \lambda_Q(L_n^-(Q))>1,\qquad
 0<\lambda_Q(L_n^+(Q))<1,
\]
and each cycle converges to the whole graphic uniformly on its ball:
\[
 \max_{\sigma\in\{-,+\}}\ \sup_{Q\in U_n}
 d_{\rm H}(L_n^\sigma(Q),\Gamma_{b_0})\longrightarrow0.
\]
Moreover,
$\sup_{Q\in U_n}\|Q-X_{b_0,0}\|\to0$.
For the local normalization $Q\mapsto(b_Q,\mu_Q)$ constructed above,
the balls may be chosen so that $\nu(Q):=-(\mu_Q)_4>0$ throughout
each $U_n$.
\end{corollary}

\begin{proof}
Take $Q_n=X_{b_0,\mu_n}$ from Theorem~\ref{thm:main}(ii) and delete
finitely many terms so that $Q_n\in\mathcal W$ and the two cycles
have the multipliers in \eqref{eq:realizing-multipliers} and each lies
within Hausdorff distance $\eta/2$ of $\Gamma_{b_0}$. Put
\[
 d_n=\max_{\sigma\in\{-,+\}}
 d_{\rm H}(L_n^\sigma,\Gamma_{b_0})\longrightarrow0.
\]
For each fixed $n$, hyperbolic persistence gives distinct continuations
$L_n^\sigma(Q)$ in disjoint tubes. Indeed the finite minimal returns
are analytic and their displacement derivatives are
$\lambda-1\ne0$, so the implicit-function theorem applies. The same
flow-box clearances preserve minimality, and finite-orbit continuous
dependence gives Hausdorff continuity; the multipliers remain on their
respective sides of one.

Shrink an open coefficient ball $U_n$ about $Q_n$ so that
$U_n\subset\mathcal W\cap B(Q_n,1/n)$ and, for every $Q\in U_n$,
\[
 d_{\rm H}(L_n^\sigma(Q),L_n^\sigma)<1/n,\qquad
 d_{\rm H}(L_n^\sigma(Q),\Gamma_{b_0})<\eta
 \quad(\sigma\in\{-,+\}).
\]
The two hyperbolic cycles each have multiplicity one.
Corollary~\ref{cor:quadratic-space} excludes a third counted isolated
cycle, proving the asserted exact count. The triangle inequality gives
\[
 \max_{\sigma\in\{-,+\}}\ \sup_{Q\in U_n}
 d_{\rm H}(L_n^\sigma(Q),\Gamma_{b_0})\le d_n+1/n\longrightarrow0,
\]
and
$\sup_{Q\in U_n}\|Q-X_{b_0,0}\|
 \le\|Q_n-X_{b_0,0}\|+1/n\to0$.

The local inverse defining the normalization gives
$\mu_{Q_n}=\mu_n$, whose fourth component is $-\nu_n<0$.
By continuity, a further shrinking of $U_n$ ensures
$(\mu_Q)_4<-\nu_n/2<0$ throughout the ball.
The radii are chosen separately for each $n$, so finite flight time
for each realizing cycle suffices even if their periods increase.
\end{proof}
